% SPDX-License-Identifier: GPL-3.0-or-later
% Copyright (C) 2026 Alessandro Greco
\chapter{Приватная сеть Tor: Chutney и Shadow}\label{cap:19-rete-privata}
\aggiornatoal{23 сентября 2026 г.}{tor 0.4.9.13; Chutney commit \texttt{6cc1588}; Shadow 3.3.0}

\begin{obiettivi}
К концу этой главы вы сможете:
\begin{itemize}[nosep]
  \item объяснять, почему эксперименты с Tor следует проводить в приватной сети, а не в публичной;
  \item с помощью двух отдельных команд собирать и запускать лабораторную сеть, используемую во всей книге;
  \item определять параметры конфигурации, отличающие тестовую сеть от реальной;
  \item ориентироваться в файлах состояния узла, включая ключи, документы консенсуса и журналы;
  \item исследовать клиентские цепи с помощью Stem и передавать через сеть обычный веб-трафик и трафик onion-сервисов;
  \item понимать, когда вместо эмулируемой сети Chutney следует использовать симулятор вроде Shadow.
\end{itemize}
\end{obiettivi}

Эта глава находится в Части~III, однако прочитать её полезно заранее: каждая лабораторная работа книги, начиная с Лабораторной работы~1 в \cref{cap:07-canali}, выполняется в сети, которую мы строим здесь. Читатели, последовательно изучающие Часть~II, могут сначала прочитать \cref{sec:rp-perche,sec:rp-avvio}, а к остальному материалу вернуться позднее.

% ---------------------------------------------------------------------------
\section{Почему нужна приватная сеть?}\label{sec:rp-perche}

Публичная сеть Tor состоит из тысяч реле, управляемых добровольцами, и ежедневно используется миллионами людей, многие из которых полагаются на неё ради собственной безопасности. Эксперименты в этой сети создают две проблемы: этическую и практическую.

\paragraph{Этическая проблема.} Любой эксперимент, который измеряет, изменяет или даже просто наблюдает трафик реальной сети, затрагивает пользователей, не дававших на это согласия и потенциально способных пострадать. Даже на первый взгляд безобидный эксперимент, например реле, записывающее больше информации, чем необходимо, может создать данные, которые при раскрытии позволят деанонимизировать человека. Поэтому Tor Project поддерживает \emph{Tor Research Safety Board}, рассматривающий экспериментальные планы исследователей, и рекомендует общий принцип: если на исследовательский вопрос можно ответить, не вовлекая реальных пользователей, следует поступать именно так. К этим принципам мы вернёмся в \cref{cap:29-ricerca}.

\paragraph{Практическая проблема.} Мы почти не контролируем публичную сеть: невозможно заранее знать, какие реле будут выбраны, цепи меняются, консенсус обновляется каждый час, а эксперимент нельзя повторить в абсолютно одинаковых условиях. В приватной сети, напротив, можно выбирать число узлов, читать подробные журналы, останавливать и перезапускать всю сеть за несколько минут и повторять эксперименты, пока результаты не станут понятны.

\begin{attenzione}[Правило лабораторных работ]
Каждое практическое упражнение этой книги, которое отправляет вручную сформированные ячейки, изменяет поведение узла или перехватывает трафик, должно выполняться \textbf{исключительно} в приватной сети. Учебный код \file{torlab} не обеспечивает анонимность и никогда не должен использоваться в публичной сети.
\end{attenzione}

% ---------------------------------------------------------------------------
\section{Эмуляция или симуляция?}\label{sec:rp-strumenti}

Существует два способа построить лабораторную сеть Tor, и оба используются Tor Project.

\begin{description}
  \item[Эмуляция] На одном компьютере запускается множество настоящих процессов \texttt{tor}, взаимодействующих через интерфейс обратной петли (loopback). Время реальное, код настоящий, а сетевое взаимодействие обеспечивает операционная система. Именно так работает \textbf{Chutney}. Chutney — инструмент Tor Project для интеграционного тестирования tor и Arti.
  \item[Симуляция] Реальные программы запускаются внутри дискретно-событийного симулятора, который перехватывает системные вызовы и предоставляет виртуальную сеть и виртуальные часы. Так работает \textbf{Shadow}~\cite{jansen2012shadow}, разработанный для масштабных исследований Tor.
\end{description}

\begin{table}[htbp]
\centering\small
\begin{tabularx}{\textwidth}{@{}l>{\raggedright\arraybackslash}X>{\raggedright\arraybackslash}X@{}}
\toprule
 & Chutney (эмуляция) & Shadow (симуляция) \\
\midrule
Типичный масштаб & От нескольких узлов до нескольких десятков & От сотен до тысяч узлов \\
Время & Реальное & Виртуальное и детерминированное \\
Сеть & Loopback без задержек и потерь & Топология с настраиваемыми задержками и пропускной способностью \\
Воспроизводимость & Ограничена планировщиком операционной системы & Полная при одинаковом случайном зерне \\
Типичное применение & Разработка, тестирование, обучение & Исследования производительности и атак \\
В этой книге & Все лабораторные работы & \cref{sec:rp-shadow} и Часть~V \\
\bottomrule
\end{tabularx}
\caption{Сравнение Chutney и Shadow.}
\label{tab:chutney-shadow}
\end{table}

Для изучения протокола естественным выбором является Chutney. Он лёгкий, запускается за несколько минут и позволяет напрямую взаимодействовать с работающими узлами. Поэтому основная часть главы посвящена Chutney, а Shadow вводится в конце.

% ---------------------------------------------------------------------------
\section{Лабораторная среда книги}\label{sec:rp-avvio}

Чтобы исключить различия между операционными системами, вся лабораторная среда упакована в Docker-образ, описанный в \file{labs/ru/лаборатория/Dockerfile} репозитория. Образ:
\begin{itemize}
  \item собирает tor \textbf{из официального исходного кода} в эталонной для книги версии 0.4.9.13, используя \texttt{--enable-gpl} для включения модуля доказательства выполнения работы (PoW) для onion-сервисов (\cref{cap:15-dos});
  \item устанавливает Chutney на зафиксированной ревизии, а также Stem и инструменты перехвата трафика (\cmd{tcpdump}, \cmd{tshark}), используемые в \cref{cap:20-traffico}.
\end{itemize}
Фиксация версий принципиальна: книга, описывающая сообщения протокола байт за байтом, должна гарантировать, что читатель увидит те же байты.

\begin{approfondimento}[Когда выходит новая версия tor]
Результаты, приведённые в этой главе и в Части~II, получены в лабораторной сети с tor 0.4.9.13 --- эталонной версией книги. Некоторые из них были записаны на предыдущей версии 0.4.9.12, выпущенной всего несколькими днями ранее, а затем воспроизведены: скрипт \file{scripts/verify\_labs.py}, который заново выполняет все лабораторные работы в живой сети и проверяет их результаты, даёт одинаковые итоги на обеих версиях (сейчас 53 проверки плюс 14 для лабораторной работы с мостом). Скрипт выполняет лабораторные работы исходного итальянского издания из \file{labs/it}, переводом кода которых является код этого издания; в листингах имена и сообщения соответствуют коду этого издания. Меняются только элементы, зависящие от конкретного запуска: ключи, отпечатки, адреса \texttt{.onion}, временные метки и размеры сертификатов, которые различаются даже между реле. Тот же скрипт позволяет вам проверить на используемой версии, что описанное в книге по-прежнему верно.
\end{approfondimento}

\subsection{Требования}

Потребуется Linux с Docker и примерно \SI{1}{\giga\byte} свободного дискового пространства. Docker Desktop работает в macOS и Windows, но используемый здесь режим сети хоста (host networking) ведёт себя иначе; в таком случае выполняйте эксперименты внутри контейнера (\cref{sec:rp-shell}). На хосте сеть использует порты, начинающиеся с 5100, 7100, 8000 и 9000. Убедитесь, что они свободны.

\subsection{Сборка образа и запуск сети}

Из \textbf{корня репозитория} (не из каталога \file{labs/ru/лаборатория}):
\begin{lstlisting}[language=bash,numbers=none]
$ make -C labs/ru/лаборатория image        # собрать tor и подготовить Chutney (один раз)
$ make -C labs/ru/лаборатория start        # настроить и запустить сеть; дождаться первоначальной загрузки
\end{lstlisting}

Сборка образа занимает несколько минут. На машине, использованной при написании этой главы, запуск сети занял примерно две с половиной минуты и завершился следующим сокращённым отчётом:
\begin{lstlisting}[numbers=none]
  test006e:
    bootstrap progress: 100 (SUCCESS), done, Done
    desc publishing: 100 (SUCCESS), all nodes, DESC MD MD_CONS NS_CONS
  test008c:
    bootstrap progress: 100 (SUCCESS), done, Done
  test009h:
    bootstrap progress: 100 (SUCCESS), done, Done
    oniondesc publishing: 200 (ONIONDESC_PUBLISHED), hidden service v3
\end{lstlisting}

Каждый узел завершил первоначальную загрузку, каждое реле видит свой дескриптор опубликованным в консенсусе, а onion-сервис опубликовал собственный дескриптор. Затем команда \cmd{make -C labs/ru/лаборатория verify} запускает тест Chutney, передавая данные как через выходное реле, так и через onion-сервис, и проверяет целостность:
\begin{lstlisting}[numbers=none]
  not-done:0 successes:4 failures:0
Transmission: Success
\end{lstlisting}

\file{make -C labs/ru/лаборатория status} показывает состояние узлов. \file{make -C labs/ru/лаборатория shell} открывает оболочку внутри контейнера, \file{make -C labs/ru/лаборатория stop} останавливает сеть, а \cmd{make -C labs/ru/лаборатория clean} также удаляет созданные данные.

% ---------------------------------------------------------------------------
\section{Топология сети}\label{sec:rp-topologia}

Сеть книги описывается коротким Python-файлом \file{labs/ru/лаборатория/chutney/networks/book}, написанным на языке конфигурации Chutney:
\begin{lstlisting}[language=Python]
Authority = Node(tag="a", authority=1, relay=1)
Middle    = Node(tag="r", relay=1)
Exit      = Node(tag="e", relay=1, exit=1)
Client    = Node(tag="c", client=1, launch_phase=2)
Onion     = Node(tag="h", hs=1, launch_phase=2)

NODES = (Authority.getN(3) + Middle.getN(3) + Exit.getN(2)
         + Client.getN(1) + Onion.getN(1))
ConfigureNodes(NODES)
\end{lstlisting}

Каждый \texttt{Node} — шаблон, а \texttt{getN} создаёт несколько экземпляров. Клиент и onion-сервис запускаются во второй фазе (\file{launch\_phase=2}), после того как авторитетные серверы сформируют консенсус. Chutney нумерует узлы по порядку и формирует имя каждого из номера и тега: \texttt{test000a}, \ldots, \texttt{test009h}. Номера узлов также определяют назначение портов, как показано на \cref{fig:rete-privata}.

\begin{figure}[htbp]
\centering
% SPDX-License-Identifier: GPL-3.0-or-later
% Copyright (C) 2026 Alessandro Greco
%% Альтернативное описание (EPUB).
\altfig{Топология тестовой сети Chutney на интерфейсе 127.0.0.1: три авторитетных сервера каталогов 000a–002a, одновременно являющиеся реле, с ORPort 5100–5102 и DirPort 7100–7102; три промежуточных реле 003r–005r, два выходных реле 006e и 007e, клиент 008c с SOCKS 9008 и onion-сервис 009h с виртуальным портом 5858, перенаправленным на 127.0.0.1:4747. В лаборатории серверы каталогов голосуют каждые 20 секунд. Пример HTTP-цепи идёт от клиента через 003r и 005r к выходному реле 006e и тестовому серверу 127.0.0.1:8080; специальная политика тестовых выходных реле разрешает локальные адреса. Onion-сервис не использует выходное реле. Схема показывает лишь часть портов, остальные приведены в таблице главы. Использование loopback не означает сетевую изоляцию от хоста.}
% Приватная лабораторная сеть: упрощённая топология
\begin{tikzpicture}[font=\sffamily\scriptsize,
  aut/.style={relay,fill=torviola!15,draw=torviola},
  ex/.style={relay,fill=torverde!30}]
  \foreach \i/\x in {0/0,1/2,2/4} \node[aut] (a\i) at (\x,3.2) {00\i a};
  \node[above=1pt of a1,font=\sffamily\scriptsize,align=center] {серверы каталогов\\ORPort 5100--5102, DirPort 7100--7102};
  \node[font=\sffamily\scriptsize,text=torgrigio,align=center] at (6.6,3.2) {голосование каждые 20 с\\подписанный консенсус};
  \foreach \i/\x in {3/0,4/2,5/4} \node[relay] (r\i) at (\x,1.2) {00\i r};
  \node[ex] (e6) at (6.6,1.2) {006e};
  \node[ex] (e7) at (6.6,-0.6) {007e};
  \node[client] (c) at (-2.4,1.2) {008c клиент};
  \node[font=\ttfamily\scriptsize,below=1pt of c] {SOCKS 9008};
  \node[client,fill=torverde!10,draw=torverde!60!black] (h) at (2,-1.2) {009h onion-сервис};
  \node[font=\ttfamily\scriptsize,below=1pt of h] {5858 $\to$ 127.0.0.1:4747};
  \node[server] (w) at (9.4,1.2) {\texttt{127.0.0.1:8080}};
  \begin{scope}[on background layer]
    \foreach \n in {r3,r4,r5,e6,e7} \draw[dotted,torgrigio] (\n) -- (a1);
  \end{scope}
  \draw[cellaf] (c) -- (r3);
  \draw[cellaf] (r3) to[bend right=40] (r5);
  \draw[cellaf] (r5) -- (e6);
  \draw[cellaf] (e6) -- (w);
  \node[font=\sffamily\scriptsize,text=torviola] at (2,-0.05) {пример HTTP-цепи};
  \node[font=\sffamily\tiny,text=torgrigio] at (4,-2.4) {onion: без exit};
\end{tikzpicture}

\caption{Топология Chutney в лаборатории на \texttt{127.0.0.1}: три авторитетных сервера каталогов (также реле), три промежуточных реле, два выходных реле, клиент и onion-сервис. ORPort $5100+n$ используется только у серверов каталогов и реле, включая выходные; ControlPort $8000+n$ имеется у каждого узла. Серверы каталогов также имеют DirPort $7100+n$, клиент — SocksPort $9008$, а onion-сервис — виртуальный порт $5858$ с перенаправлением на \texttt{127.0.0.1:4747}. Пример HTTP-цепи через выходное реле не относится к onion-сервису: тот не использует exit. Loopback не является отдельной сетью контейнера.}
\label{fig:rete-privata}
\end{figure}

\begin{center}\small
\setlength{\tabcolsep}{3.5pt}
\begin{tabular}{lllll}
\toprule
Узел & Роль & ORPort & ControlPort & Дополнительно \\
\midrule
\texttt{test000a}--\texttt{002a} & авторитетный сервер каталогов & 5100--5102 & 8000--8002 & DirPort 7100--7102 \\
\texttt{test003r}--\texttt{005r} & промежуточное реле & 5103--5105 & 8003--8005 & \\
\texttt{test006e}--\texttt{007e} & выходное реле & 5106--5107 & 8006--8007 & \\
\texttt{test008c} & клиент & --- & 8008 & SocksPort 9008 \\
\texttt{test009h} & onion-сервис & --- & 8009 & виртуальный порт 5858 \\
\bottomrule
\end{tabular}
\end{center}

\begin{approfondimento}[Почему три авторитетных сервера каталогов?]
Консенсус должен быть подписан большинством авторитетных серверов каталогов (\cref{cap:12-directory}). С одним таким сервером мы не смогли бы наблюдать голосование; с тремя можно увидеть весь процесс и в упражнении проверить, что происходит после остановки одного из них. В публичной сети их девять плюс отдельный авторитетный сервер мостов. Список встроен непосредственно в tor (\file{src/app/config/auth\_dirs.inc}).
\end{approfondimento}

% ---------------------------------------------------------------------------
\section{Тестовая сеть — не реальная сеть}\label{sec:rp-opzioni}

Chutney генерирует отдельный \file{torrc} для каждого узла. Эти файлы показывают, что именно необходимо изменить для работы миниатюрной сети Tor и, следовательно, какие предположения о публичной сети здесь неприменимы. Ниже приведены наиболее важные параметры из torrc одного авторитетного сервера (\file{net/nodes/000a/torrc}):
\begin{lstlisting}[language=torrc,numbers=none]
TestingTorNetwork 1
AuthoritativeDirectory 1
V3AuthoritativeDirectory 1
AssumeReachable 1
V3AuthVotingInterval 20 seconds
V3AuthVoteDelay 4
V3AuthDistDelay 4
TestingDirAuthVoteGuard *
TestingDirAuthVoteHSDir *
\end{lstlisting}

\begin{description}
  \item[\opt{TestingTorNetwork}] Включает тестовый режим. Его можно использовать только вместе с явным списком авторитетных серверов каталогов (\opt{DirAuthority}), что не позволяет тестовым узлам случайно присоединиться к публичной сети. Кроме того, этот режим изменяет многие значения по умолчанию ради ускорения работы.
  \item[\opt{V3AuthVotingInterval}] В реальной сети авторитетные серверы голосуют за новый консенсус каждый час. Здесь голосование проходит каждые 20 секунд, что видно и в самом консенсусе:
\begin{lstlisting}[numbers=none]
valid-after 2026-09-23 14:37:40
fresh-until 2026-09-23 14:38:00
valid-until 2026-09-23 14:38:20
\end{lstlisting}
  \item[\opt{AssumeReachable}] В обычной сети реле сначала проверяет, доступен ли он извне, и лишь затем публикует себя. На интерфейсе обратной петли такая проверка бессмысленна, поэтому пропускается.
  \item[\opt{TestingDirAuthVoteGuard} \texttt{*}] В реальной сети реле получает флаг Guard только после недель стабильной работы и при достаточной пропускной способности. Здесь авторитетные серверы присваивают этот флаг всем реле. Поэтому \textbf{каждое реле нашего консенсуса может быть входным узлом (guard)}, чего в публичной сети нет.
\end{description}

Файлы torrc выходных реле также содержат:
\begin{lstlisting}[language=torrc,numbers=none]
ExitPolicyRejectPrivate 0
ExitPolicy accept private:*
\end{lstlisting}
По умолчанию выходные реле отклоняют соединения к приватным адресам и адресам обратной петли, чтобы не открывать доступ к собственной локальной сети. В этой лаборатории \texttt{127.0.0.1} — единственный доступный \enquote{Интернет}, поэтому политика должна быть обратной.

\begin{attenzione}[Не копируйте эти параметры в другие среды]
Опции \texttt{Testing\ldots}, \opt{AssumeReachable} и особенно \texttt{ExitPolicy accept private:*} предназначены исключительно для тестовых сетей. На публичном реле такие настройки были бы неверными или опасными для сети хоста.
\end{attenzione}

% ---------------------------------------------------------------------------
\section{Внутри каталога узла}\label{sec:rp-file}

У каждого узла есть собственный каталог данных в \file{labs/ru/лаборатория/net/nodes/NNNx/}. Среди прочего, каталог авторитетного сервера содержит:

\begin{center}\small
\begin{tabularx}{\textwidth}{@{}l>{\raggedright\arraybackslash}X@{}}
\toprule
Файл или каталог & Содержимое \\
\midrule
\file{torrc} & конфигурация, сгенерированная Chutney \\
\file{keys/} & идентификационные и onion-ключи узла (\cref{cap:05-crittografia,cap:06-architettura}) \\
\file{fingerprint}, \file{fingerprint-ed25519} & отпечатки идентичности RSA и Ed25519 \\
\file{cached-consensus}, \file{cached-microdesc-consensus} & оба варианта консенсуса (\cref{cap:12-directory}) \\
\file{cached-microdescs.new} & загруженные микродескрипторы \\
\file{v3-status-votes} & голоса авторитетных серверов (только для них) \\
\file{sr-state} & состояние протокола общей случайности (только для авторитетных серверов) \\
\file{notice.log}, \file{info.log} & журналы уровней \emph{notice} и \emph{info} \\
\file{control\_auth\_cookie} & cookie для аутентификации через порт управления \\
\bottomrule
\end{tabularx}
\end{center}

Каталог onion-сервиса дополнительно содержит \file{hidden\_service/} с его ключами и файл \file{hostname} с адресом \texttt{.onion}, например:
\begin{lstlisting}[numbers=none]
rpns2r36zn5u46wctlgnvtwjotxaf7iflsc3k5eas3lqoxeesv4zwyad.onion
\end{lstlisting}
Адрес меняется при каждой новой конфигурации, поскольку генерируется новый ключ (\cref{cap:14-onion-services}).

Консенсус — обычный человекочитаемый текстовый файл. Если извлечь строки \texttt{r} (идентификаторы реле и адреса) и \texttt{s} (флаги), получим, например, такие два реле:
\begin{lstlisting}[numbers=none]
r test000a H+9DDHx4rEo6I6Jdn8gIyKyi3jk ... 127.0.0.1 5100 7100
s Authority Fast Guard HSDir Running Stable V2Dir Valid
r test006e dEPHoFFHamQyy6pRzlsLB9/ilzU ... 127.0.0.1 5106 0
s Exit Fast Guard HSDir Running Stable V2Dir Valid
\end{lstlisting}
Как и ожидалось, все реле имеют флаг Guard, а Exit — только два выходных реле. Поля подробно разбираются в \cref{cap:12-directory}.

% ---------------------------------------------------------------------------
\section{Использование сети}\label{sec:rp-uso}

\subsection{Доступ через приватную сеть}

Оба HTTP-примера выполняйте \emph{внутри лабораторного контейнера}, в каталоге \file{/lab}, открыв его из отдельного терминала командой \file{make -C labs/ru/лаборатория shell}. Пустая строка в параметре curl \texttt{--noproxy} принуждает использовать SOCKS-прокси даже для \texttt{127.0.0.1}, если задана переменная \texttt{NO\_PROXY}. Сервер на порту 4747 использует тот же \file{index.html}, который создан в первом примере.

Клиент \texttt{test008c} предоставляет SOCKS на порту 9008. Его может использовать любая программа с поддержкой SOCKS5. Запустим небольшой локальный веб-сервер и обратимся к нему через Tor:
\begin{lstlisting}[language=bash,numbers=none]
$ echo "Привет из приватной сети" > index.html
$ python3 -m http.server 8080 --bind 127.0.0.1 &
$ curl --noproxy "" --socks5-hostname 127.0.0.1:9008 http://127.0.0.1:8080/index.html
Привет из приватной сети
\end{lstlisting}
Запрос прошёл по цепи из трёх реле, после чего одно из выходных реле подключилось к серверу. Использование \texttt{--socks5-hostname} вместо \texttt{--socks5} заставляет exit-реле разрешать имена назначения вместо локальной системы пользователя. В реальной сети это принципиально для предотвращения утечки посещаемых адресов локальному DNS-резолверу (\cref{cap:11-stream}).

\subsection{Доступ к onion-сервису}

Chutney настраивает onion-сервис так:
\begin{lstlisting}[language=torrc,numbers=none]
HiddenServicePort 5858 127.0.0.1:4747
\end{lstlisting}
Это сопоставляет виртуальный порт 5858 адреса \texttt{.onion} с локальным портом 4747. Остаётся лишь запустить сервер:
\begin{lstlisting}[language=bash,numbers=none]
$ python3 -m http.server 4747 --bind 127.0.0.1 &
$ curl --noproxy "" --socks5-hostname 127.0.0.1:9008 \
    http://$(cat net/nodes/009h/hidden_service/hostname):5858/index.html
Привет из приватной сети
\end{lstlisting}
На этот раз выходное реле не участвует. Клиент и сервис встречаются в пункте встречи (rendezvous) внутри сети (\cref{cap:14-onion-services}).

\subsection{Наблюдение за цепями с помощью Stem}\label{sec:rp-stem}

Порт управления позволяет исследовать работающий узел. С помощью Stem (\cref{cap:18-stem}), включённого в образ, выведем цепи клиента:
\begin{lstlisting}[language=Python]
from stem.control import Controller

with Controller.from_port(port=8008) as c:
    c.authenticate()                    # cookie используется автоматически
    print(c.get_version())
    for circ in c.get_circuits():
        print(circ.id, circ.status, circ.purpose,
              [name for _, name in circ.path])
\end{lstlisting}
Фрагмент вывода после \cmd{make -C labs/ru/лаборатория verify}:
\begin{lstlisting}[numbers=none]
0.4.9.13 (git-3c575400909efe65)
10 BUILT GENERAL ['test004r', 'test005r', 'test007e']
124 BUILT CONFLUX_LINKED ['test006e', 'test004r', 'test007e']
114 BUILT CONFLUX_UNLINKED ['test007e', 'test004r', 'test006e']
128 BUILT CIRCUIT_PADDING ['test004r', 'test001a', 'test007e', 'test006e']
163 BUILT HS_CLIENT_REND ['test004r', 'test001a', 'test006e']
166 BUILT HS_CLIENT_HSDIR ['test004r', 'test006e', 'test007e', 'test003r']
181 BUILT HS_VANGUARDS ['test004r', 'test007e', 'test006e']
\end{lstlisting}

Даже в столь небольшой сети список предваряет значительную часть книги:
\begin{itemize}
  \item Почти все цепи начинаются с \texttt{test004r}: это \textbf{входное реле (guard)} клиента, выбранный один раз и затем многократно используемый (\cref{cap:13-percorso}). Исключение составляют ветви туннеля Conflux, начинающиеся со второго входного реле (\cref{cap:18-stem}).
  \item Трёхузловые цепи \texttt{GENERAL} переносят обычный трафик; одношаговые цепи, не показанные здесь, используются для загрузки документов каталога.
  \item Цепи \file{CONFLUX\_LINKED} появляются парами и совместно переносят один и тот же трафик: это механизм Conflux (\proposta{329}, \cref{cap:31-futuro}).
  \item \file{HS\_VANGUARDS}, \file{HS\_CLIENT\_HSDIR} и \file{HS\_CLIENT\_REND} обозначают цепи, участвующие в доступе к onion-сервису (\cref{cap:14-onion-services}).
  \item Согласно спецификации протокола управления, \file{CIRCUIT\_PADDING} обозначает цепь, сохраняемую после окончания реального использования, чтобы скрыть точный момент завершения активности. Поэтому её закрытие не обязательно показывает, когда завершилась пользовательская активность.
\end{itemize}

\subsection{Работа внутри контейнера}\label{sec:rp-shell}

Контейнер использует сетевой режим хоста (host networking), поэтому все порты узлов доступны непосредственно с хостовой системы. Если этот режим недоступен, например в Docker Desktop, \file{make -C labs/ru/лаборатория shell} открывает оболочку внутри контейнера, где установлены tor, Chutney, Stem, Python и инструменты перехвата; лабораторные файлы доступны в \file{/lab}.

% ---------------------------------------------------------------------------
\section{Первый контакт с реле}\label{sec:rp-contatto}

Следующий фрагмент Python импортирует учебный пакет \file{torlab} из \file{labs/ru/код/}; контейнер Chutney монтирует только \file{labs/ru/лаборатория/} в \file{/lab}, поэтому пакет \emph{не доступен в контейнере автоматически}. В Linux с режимом сети host сохраните фрагмент в файл и запустите из корня репозитория с \texttt{PYTHONPATH=labs/ru/код python3 script.py} либо явно предоставьте пакет контейнеру. Пример предназначен исключительно для приватной сети.

Вернёмся к \cref{cap:08-celle}. Используя \file{labs/ru/код/torlab/cells.py} из Лабораторной работы~2, можно открыть TLS-соединение с ORPort лабораторного реле, отправить ячейку \cella{VERSIONS} и декодировать ответ. Лабораторная работа~1 в \cref{cap:07-canali} содержит полный код с проверкой сертификатов. Здесь мы лишь наблюдаем начальный обмен:
\begin{lstlisting}[language=Python]
import socket, ssl
from torlab.cells import *

ctx = ssl.SSLContext(ssl.PROTOCOL_TLS_CLIENT)
ctx.check_hostname = False
ctx.verify_mode = ssl.CERT_NONE   # только приватная сеть; сертификаты TLS НЕ проверяются
s = ctx.wrap_socket(socket.create_connection(("127.0.0.1", 5103)))
s.sendall(encode(versions_cell([3, 4, 5]), 0))

buf, version, received = b"", 0, 0
while received < 4:
    chunk = s.recv(4096)
    if not chunk:
        raise ConnectionError("TLS channel closed early")
    buf += chunk
    while (r := decode(buf, version)) is not None:
        cell, buf = r
        received += 1
        print(describe(cell, version))
        if cell.command == CellCommand.VERSIONS:
            version = negotiate_version([3, 4, 5], cell.body)
s.close()
\end{lstlisting}

\textbf{Предупреждение:} эта краткая демонстрация намеренно отключает проверку сертификатов TLS через \texttt{ssl.CERT\_NONE} и \emph{не} проверяет подлинность реле по ячейке CERTS. Не выполняйте её на публичных реле. Полная процедура проверки приведена в лабораторной работе~1.

Следующий вывод взят из предыдущего эталонного запуска: идентификаторы цепей, отпечатки, onion-адреса, временные метки и размеры сертификатов зависят от конкретного эксперимента.

Реле \texttt{test003r} отвечает:
\begin{lstlisting}[numbers=none]
CircID=0x0000 VERSIONS тело=6 байт
CircID=0x00000000 CERTS тело=1451 байт
CircID=0x00000000 AUTH_CHALLENGE тело=36 байт
CircID=0x00000000 NETINFO тело=509 байт
\end{lstlisting}

Это ровно та последовательность, которая описана в \cref{sec:sessione}: сначала \cella{VERSIONS} с двухбайтовым CircID; после согласования версии 5 последующие ячейки используют четырёхбайтовый CircID. \cella{CERTS} и \cella{AUTH\_CHALLENGE} имеют переменную длину и сохраняют точные тела сообщений; \cella{NETINFO} имеет фиксированную длину и несёт все 509 байт. Обратите внимание, что \texttt{version = ...} присваивается сразу после \cella{VERSIONS}: декодирование следующих ячеек зависит от только что согласованной версии.

% ---------------------------------------------------------------------------
\section{Shadow: сеть внутри симулятора}\label{sec:rp-shadow}

Chutney недостаточно для вопросов о поведении сети \emph{в целом}. Как другой алгоритм управления перегрузкой повлиял бы на задержку? Какой объём трафика понадобился бы атакующему для определённой атаки? Для ответа нужны тысячи узлов, реалистичные задержки между континентами и, прежде всего, воспроизводимые эксперименты.

Shadow, первоначально разработанный Робом Янсеном и Николасом Хоппером~\cite{jansen2012shadow}, решает эти задачи. Современные версии запускают реальные программы, включая tor, как обычные Linux-процессы, одновременно перехватывая сетевые и связанные со временем системные вызовы. Виртуальный граф задаёт пользовательские задержки и пропускную способность, а время продвигается только тогда, когда это решает симулятор. Поэтому один симулированный час может занимать больше или меньше одного реального часа в зависимости от вычислительной мощности, а повтор эксперимента с тем же случайным зерном даёт тот же результат.

Для построения реалистичных Tor-сетей в Shadow исследователи используют \textbf{tornettools}. На основе публичных данных Tor Metrics (\cref{cap:24-metrics}) он создаёт уменьшенную сеть с распределениями пропускной способности, флагов и ролей, соответствующими реальной сети.

Мы используем Shadow в управляемом эксперименте в \cref{cap:29-ricerca}. Пока достаточно понимать, когда именно симулятор является подходящим инструментом.

% ---------------------------------------------------------------------------
\section*{Итоги}
\begin{itemize}
  \item Проводите эксперименты с Tor в приватных сетях. В публичной сети они могут затрагивать пользователей без их согласия и не обеспечивают воспроизводимость.
  \item Chutney эмулирует Tor-сеть настоящими процессами \texttt{tor} на интерфейсе обратной петли. Shadow запускает настоящее программное обеспечение внутри дискретно-событийного симулятора, подходящего для больших сетей и повторяемых экспериментов.
  \item Лабораторная среда книги (\file{labs/ru/лаборатория/}) собирает tor 0.4.9.13 из исходного кода. \cmd{make -C labs/ru/лаборатория start} запускает сеть из десяти узлов: три авторитетных сервера, три промежуточных реле, два выходных реле, клиент и onion-сервис.
  \item Узел $n$ использует ORPort $5100+n$, ControlPort $8000+n$ и SocksPort $9000+n$.
  \item Тестовые настройки (\opt{TestingTorNetwork}, голосование каждые 20 секунд, назначение Guard всем реле и разрешение приватных адресов для выходных реле) позволяют миниатюрной сети работать, но отличают её от реальной. Результаты измерений требуют осторожности при экстраполяции.
  \item SocksPort 9008 передаёт прикладной трафик через сеть; Stem наблюдает цепи; \file{torlab} напрямую взаимодействует с реле.
\end{itemize}

% ---------------------------------------------------------------------------
\section*{Упражнения}

\begin{esercizio}[1]\label{es:19-porte}
Не читая torrc, назовите ORPort и ControlPort узла \texttt{test006e}. Есть ли ORPort у \texttt{test009h}? Почему?
\end{esercizio}

\begin{esercizio}[1]\label{es:19-guard}
Каждое реле приватной сети имеет флаг Guard. Какая опция вызывает это поведение и почему в реальной сети всё иначе? Как это влияет на лабораторные измерения выбора маршрута?
\end{esercizio}

\begin{esercizio}[1]\label{es:19-socks}
Повторите запрос \texttt{curl} из \cref{sec:rp-uso}, используя \texttt{--socks5} вместо \texttt{--socks5-hostname} и \texttt{localhost} вместо \texttt{127.0.0.1}. Что изменится в данных, получаемых exit-реле? Почему это важно в публичной сети?
\end{esercizio}

\begin{esercizio}[2]\label{es:19-consenso}
Несколько минут наблюдайте поле \texttt{valid-after} консенсуса клиента (\cmd{grep valid-after net/nodes/008c/cached-microdesc-consensus}). Как часто оно меняется? Вычислите, сколько консенсусов создаётся за сутки в лабораторной и реальной сетях.
\end{esercizio}

\begin{esercizio}[2]\label{es:19-autorita}
Остановите одну из трёх authorities (\file{make -C labs/ru/лаборатория shell}, затем \cmd{kill} соответствующего процесса, используя его файл \file{pid}). Продолжает ли сеть создавать консенсусы? Что происходит после остановки двух authorities? Объясните результат; основная причина рассмотрена в \cref{cap:12-directory}.
\end{esercizio}

\begin{esercizio}[2]\label{es:19-stem}
Измените сценарий Stem из \cref{sec:rp-stem}, чтобы он выводил ORPort всех реле в каждой трёхузловой цепи, получая их из консенсуса через \file{get\_network\_status}.
\end{esercizio}

\begin{esercizio}[2]\label{es:19-rete}
Создайте \file{labs/ru/лаборатория/chutney/networks/book-large} с пятью авторитетными серверами, десятью промежуточными и четырьмя выходными реле, затем запустите его, установив \texttt{NETWORK=book-large} в команде \texttt{make}. Как время первоначальной загрузки соотносится со временем запуска сети книги?
\end{esercizio}

\begin{esercizio}[3]\label{es:19-log}
Найдите в \file{info.log} клиента сообщения, относящиеся к построению цепи, начиная с выбора узлов и заканчивая завершением. Сопоставьте каждое сообщение с ячейкой на \cref{fig:sessione}.
\end{esercizio}

\begin{esercizio}[3]\label{es:19-cattura}
Внутри контейнера перехватите трафик к порту 5103 командой \cmd{tcpdump -i lo} во время запуска сценария из \cref{sec:rp-contatto}. Видны ли ячейки? Почему нет? Что всё ещё остаётся наблюдаемым: размеры, временные интервалы? \Cref{cap:20-traffico} показывает, как исследовать глубже.
\end{esercizio}

% ---------------------------------------------------------------------------
\section*{Дополнительная литература}
\begin{itemize}
  \item Chutney, \url{https://gitlab.torproject.org/tpo/core/chutney}: README и каталог \file{networks/} содержат множество примеров сетей, включая конфигурации с мостами и Arti.
  \item R.~Jansen, N.~Hopper, \titolo{\textenglish{Shadow: Running Tor in a Box for Accurate and Efficient Experimentation}}~\cite{jansen2012shadow}.
  \item Документация Shadow и tornettools: \url{https://shadow.github.io/}.
  \item Tor Research Safety Board, \url{https://research.torproject.org/safetyboard/}: рекомендации по этическим исследованиям Tor.
\end{itemize}
